\usepackage{avant} 							%Avant als normale Schrift
\usepackage[helvet]{sfmath} 					%serifenfreie Schrift für Mathmode


\renewcommand{\familydefault}{\sfdefault}
\usepackage[onehalfspacing]{setspace}
\usepackage[utf8]{inputenc}
\usepackage[pdftex]{graphicx} 					% Einbindung von Grafiken (pdf, png, jpg)
\usepackage{float}            					% bietet Option [H] für bombenfestes Verankern
\usepackage[ngerman]{babel}   					% Silbentrennung nach der neuen deutschen Rechtschreibung
\usepackage{amstext}          					% für Klartext via \text{} in Formeln
\usepackage{amsfonts}
\usepackage{amsmath}        				    % für komplexere Formeln
\usepackage{amssymb}        					% für komplexere Formeln
\usepackage{bm}           					    % bold math, für \bm{}
\usepackage{enumerate}        					% verbessert Aufzählungen
\usepackage[bottom]{footmisc} 					% Fussnoten am Seitenende
\interfootnotelinepenalty=10000 				% Verhindert Fußnoten-Trennung über Seiten
\usepackage{array}            					% für Tabellen: bindet tabular-Umgebung ein
\usepackage{tabularx}						    % für Tabellen mit automatischer Spaltenbreite
\usepackage{longtable}        					% für Tabellen mit Seitenumbruch
\usepackage{algorithm}        					% für Algorithmen
\usepackage{algorithmic}      					% für Algorithmen
\usepackage{ntheorem}
\usepackage{theorem}
\usepackage{pdfpages}         					% für die Einbindung kompletter pdf-Seiten
\usepackage{parskip}          					% zwischen Absätzen: eine knappe Leerzeile
\usepackage[right]{eurosym}   					% Eurosymbol
\usepackage{xcolor}           					% farbiger Text
\usepackage{color}
\definecolor{OBlue}{rgb}{0.0,0.2,0.4}
\definecolor{OOrange}{rgb}{1.0,0.55,0.0}
\PassOptionsToPackage{hyphens}{url}
\usepackage[colorlinks=true,linkcolor=OBlue,citecolor=black]{hyperref}
\usepackage[backend=biber, style=numeric, citestyle=numeric-comp]{biblatex}
\DeclareFieldFormat{labelnumberwidth}{\mkbibbrackets{#1}}
\setlength{\bibitemsep}{1\baselineskip}
\usepackage{wasysym}
\usepackage{pdfpages}
\usepackage{epstopdf}
\usepackage[nice]{nicefrac}					    % für Brüche im Text
\usepackage{fancyhdr}						    % Kopf und Fußzeile ermöglichen
\usepackage[T1]{fontenc}
\newcommand{\changefont}[3]{\fontfamily{#1} \fontseries{#2} \fontshape{#3} \selectfont}
\changefont{ptm}{m}{n}
\linespread {1.5}

\sloppy
\usepackage{titlesec}
\titleformat{\chapter}{\bf\Huge}{\thechapter\quad}{0em}{}
\let\cleardoublepage\clearpage

\usepackage{verbatim}
\usepackage{tikz}
\usetikzlibrary{arrows.meta}
\usepackage{pgfplots}
\pgfplotsset{compat=1.18}

% Größenanpassungen
\setlength{\unitlength}{1cm}
\setlength{\oddsidemargin}{0.3cm}
\setlength{\evensidemargin}{0.3cm}
\setlength{\textwidth}{16cm}
\setlength{\topmargin}{-1.2cm}
\setlength{\textheight}{23cm}
\columnsep 0.5cm

\makeindex
